\documentclass[11pt,a4paper]{article}
% =====================================================================
%  Shared preamble for the Part V lecture notes (lectures 19-22).
%
%  These four lectures sit outside Geron, so the notes ARE the primary
%  source and are examined on the same terms as the textbook parts.
%  Everything here is chosen to make that readable on paper: one column,
%  generous measure, and the same six-colour palette as the slides so a
%  student moving between the two is not re-learning what red means.
%
%  Build with pdflatex (twice, for the toc and the refs):
%      cd notes && latexmk -pdf lecture-19.tex
%  or  make -C notes
% =====================================================================
\usepackage[T1]{fontenc}
\usepackage[utf8]{inputenc}
\usepackage{newtxtext,newtxmath}      % Times-like: prints well, reads well
\usepackage[scaled=0.86]{inconsolata} % code, at a size that sits beside Times
\usepackage{microtype}
% newtx defines \Bbbk, \openbox, \proof and \endproof; amssymb and amsthm
% then redefine them and abort. Freeing the four names before those two
% packages load is the standard fix, and it costs nothing here: we use
% amsthm only for \newtheorem.
\let\Bbbk\relax
\usepackage{amsmath,amssymb}
\let\openbox\relax
\let\proof\relax
\let\endproof\relax
\usepackage{amsthm}
\usepackage{booktabs}
\usepackage{enumitem}
\usepackage{xcolor}
\usepackage{tcolorbox}
\usepackage{titlesec}
\usepackage{graphicx}
\usepackage[hidelinks]{hyperref}
\tcbuselibrary{skins,breakable}

% ---- the course palette, from assets/css/custom.css -------------------
\definecolor{aimlprimary}{HTML}{0B3D62}   % deep blue  - structure
\definecolor{aimlaccent} {HTML}{C0392B}   % brick red  - failures
\definecolor{aimlsuccess}{HTML}{14663A}   % green      - repairs
\definecolor{aimlmath}   {HTML}{6C3483}   % purple     - the derivation
\definecolor{aimlmuted}  {HTML}{4B5563}
\definecolor{aimlrule}   {HTML}{B0BCC7}
\definecolor{aimlsoft}   {HTML}{F4F7F9}

\usepackage[a4paper,margin=32mm,top=30mm,bottom=30mm]{geometry}
\setlength{\parskip}{0.45\baselineskip}
\setlength{\parindent}{0pt}
\linespread{1.06}

% ---- headings --------------------------------------------------------
\titleformat{\section}{\Large\bfseries\color{aimlprimary}}{\thesection}{0.7em}{}
\titleformat{\subsection}{\large\bfseries\color{aimlprimary}}{\thesubsection}{0.7em}{}
\titleformat{\subsubsection}{\bfseries\color{aimlmuted}}{}{0em}{}
\titlespacing*{\section}{0pt}{1.6\baselineskip}{0.5\baselineskip}
\titlespacing*{\subsection}{0pt}{1.2\baselineskip}{0.35\baselineskip}

% ---- boxes -----------------------------------------------------------
% One box per role, and the role is the colour. A student who has seen the
% slides already knows what each one means before reading a word of it.
\newtcolorbox{keybox}[1][]{
  breakable, enhanced, colback=aimlsoft, colframe=aimlprimary,
  boxrule=0.4pt, left=8pt, right=8pt, top=6pt, bottom=6pt,
  fonttitle=\bfseries\small, coltitle=white, colbacktitle=aimlprimary,
  attach boxed title to top left={yshift=-2mm, xshift=4mm},
  boxed title style={boxrule=0pt, arc=1pt}, #1}

\newtcolorbox{failbox}[1][]{
  breakable, enhanced, colback=white, colframe=aimlaccent,
  boxrule=0.9pt, left=8pt, right=8pt, top=6pt, bottom=6pt,
  fonttitle=\bfseries\small, coltitle=white, colbacktitle=aimlaccent,
  attach boxed title to top left={yshift=-2mm, xshift=4mm},
  boxed title style={boxrule=0pt, arc=1pt}, #1}

\newtcolorbox{derivbox}[1][]{
  breakable, enhanced, colback=white, colframe=aimlmath,
  boxrule=0.6pt, left=8pt, right=8pt, top=6pt, bottom=6pt,
  fonttitle=\bfseries\small, coltitle=white, colbacktitle=aimlmath,
  attach boxed title to top left={yshift=-2mm, xshift=4mm},
  boxed title style={boxrule=0pt, arc=1pt}, #1}

% An exercise. Deliberately the same shape as the notebook's `try` field:
% one change, and what should happen to the output.
\newtcolorbox{trybox}[1][]{
  breakable, enhanced, colback=white, colframe=aimlsuccess,
  boxrule=0.5pt, left=8pt, right=8pt, top=6pt, bottom=6pt,
  fonttitle=\bfseries\small, coltitle=white, colbacktitle=aimlsuccess,
  attach boxed title to top left={yshift=-2mm, xshift=4mm},
  boxed title style={boxrule=0pt, arc=1pt}, #1}

% ---- small semantics -------------------------------------------------
\newcommand{\scope}[1]{\textcolor{aimlmuted}{\small #1}}
\newcommand{\term}[1]{\textbf{#1}}
\newcommand{\code}[1]{\texttt{\small #1}}
\newcommand{\lecture}[1]{Lecture~#1}

\newtheorem{definition}{Definition}[section]
\newtheorem{proposition}[definition]{Proposition}

% ---- title block -----------------------------------------------------
% \notestitle{19}{Information retrieval: the lexical foundation}{SciFact (BEIR)}
\newcommand{\notestitle}[3]{%
  \thispagestyle{empty}
  \begin{flushleft}
    {\color{aimlmuted}\small\sffamily
     APPLICATIONS OF MACHINE LEARNING \quad$\cdot$\quad
     BSc Mathematics of Artificial Intelligence}\\[2mm]
    {\color{aimlrule}\rule{\textwidth}{0.8pt}}\\[4mm]
    {\color{aimlmuted}\large Lecture #1 \quad$\cdot$\quad Part V \quad$\cdot$\quad
     Extended notes}\\[3mm]
    {\Huge\bfseries\color{aimlprimary}\baselineskip=1.15\baselineskip #2\par}\vspace{4mm}
    {\color{aimlmuted}\large Dataset: #3}\\[3mm]
    {\color{aimlrule}\rule{\textwidth}{0.8pt}}
  \end{flushleft}
  \vspace{2mm}
  \begin{keybox}[title={Why these notes exist}]
    Lectures 19--22 sit outside G\'eron, the course's primary text. For those
    four lectures \emph{these notes are the primary source}, and they are
    examined on exactly the same terms as the textbook parts. Everything in
    the slide deck for this lecture is here, with the arguments written out at
    length rather than compressed onto a slide; the numbers are the ones the
    accompanying notebook prints, and every one of them is a measurement on
    the corpus named above rather than a figure quoted from a paper.
  \end{keybox}
  \vspace{1mm}
  {\small\tableofcontents}
  \vspace{2mm}
  \noindent{\color{aimlrule}\rule{\textwidth}{0.4pt}}
  \vspace{1mm}
}


\begin{document}
\notestitlefor{23}{VI}{Vision transformers and multimodal retrieval}{COCO, 200 catalogue entries}{Chapters 15--16}

\section{The first brief whose input and output differ in kind}

Part V built retrieval over text and over user behaviour, and both reduced to
the same object: two encoders and an inner product. Today the two towers take
\emph{different modalities} --- one encodes images, the other text, and a
photograph and its caption must land close together.

A retailer's catalogue team: the entries are photographs. Product photos arrive
with an identifier, a price and a picture; written descriptions are
inconsistent, often absent, and never written by the person searching.
\emph{Customers type sentences. The catalogue stores pixels.}

Keyword search needs text on both sides, and so does the semantic search of the
previous lecture --- it embeds a sentence and compares it to \emph{other
sentences}. An entry with no description is not badly ranked by a text index.
It is \emph{absent}.

\subsection{The corpus, and the second corpus}

200 catalogue entries, 32.6 MB, taken in image-id order from COCO's 2014
validation set so that everyone's catalogue is the same catalogue, each with
five independent human captions written by people who were not shown each
other's.

\begin{keybox}[title={Two different uses for five captions}]
Caption 1 is the entry's \emph{description}, when it has one --- what a text
index would see. Caption 2 is the \emph{customer's query}. Captions 3 to 5 are
held back.

Using two different people's sentences about the same picture makes this a
paraphrase problem rather than a lookup. If we used the same sentence on both
sides, a hash table would score 100\%.
\end{keybox}

The catalogue has no labels, so it cannot play the role of a known-answer test.
A second corpus does: 2{,}000 CIFAR-10 test images and ten class names --- ten
labels we actually have. CIFAR-10 is $32\times32$, which is a hard setting for
any model trained on web photographs, and we will see the cost.

\subsection{An image as a sequence}

A transformer takes a sequence and an image is a grid, so cut it up: split into
fixed patches, flatten each and project it linearly --- \emph{that projection
is Lecture 17's embedding table with patches in place of tokens} --- add a
positional encoding, because Lecture 18 established that attention has no
notion of order and a grid has two, then run the standard blocks. A
$224\times224$ image at $16\times16$ patches is $14\times14 = 196$ tokens, well
within reach of the quadratic cost, which is exactly why that patch size and
not a smaller one.

\begin{keybox}[title={What a ViT gives up}]
Lecture 12 derived what convolution assumes: locality, and translation
equivariance from weight sharing. A vision transformer assumes \emph{neither}.
Every patch can attend to every other from the first layer, and nothing tells
it that adjacent patches are related.

On small corpora a convolutional network wins, because its assumptions are true
and free. On large ones the transformer wins, because it can learn whichever
relations actually hold. The crossover is at a corpus size most people do not
have --- which is why pretrained ViTs are used far more often than they are
trained.
\end{keybox}

\section{A metric, and a baseline that needs no experiment}

The system returns a ranked list, so what matters is where the right entry
landed. Recall@$k$ is the share of queries whose answer is in the top $k$, and
we read R@1, R@5 and R@10 --- R@1 being the number the product cares about,
because it is the top of the page. Ties count \emph{against} us: a metric that
flatters a degenerate score matrix is not a metric.

Not accuracy, because there is no decision here, only an ordering --- and a
threshold on similarity would import all of Lecture 3's problems plus a new
one, that the scale of a cosine is not comparable between models. Recall@$k$ is
rank-based, so it is invariant to any monotone transformation of the score.
Mean reciprocal rank goes beside it, so that a system putting the answer at 2
is distinguishable from one putting it at 180.

\begin{derivbox}[title={The trivial baseline, computed rather than run}]
One relevant entry among $n$, placed uniformly at random:
\[ \Pr[r \le k] = \frac{k}{n}, \qquad
   \mathbb{E}[r] = \frac{n+1}{2}, \qquad
   \mathbb{E}\!\left[\tfrac{1}{r}\right] = \frac{H_n}{n}. \]
No simulation, no seed, no error bar --- this one is arithmetic.
\end{derivbox}

At $n = 200$ that is R@1 of 0.5\%, R@5 of 2.5\%, R@10 of 5.0\%, an expected
rank of 100.5 and an MRR of 0.029. These scale with the catalogue: on 5{,}000
candidates R@1 would be 0.02\%. \emph{Quote the size or quote nothing.}

\section{Two encoders, trained apart}

Encode the images with a ViT trained on ImageNet labels --- no text anywhere in
its training --- and the queries with the sentence encoder from the previous
lecture. Then take the cosine.

\begin{failbox}[title={It does not even typecheck}]
The image vectors live in 768 dimensions and the sentence vectors in 384.
There is no cosine between them, because there is no inner product between
them.

Reviewer question 3, asked by the interpreter instead of by you.
\end{failbox}

Make the dimensions agree three ways, so nobody can blame the fix: a random
Gaussian projection (Lecture 8), keeping the first 384 coordinates, or
zero-padding the text vectors. Lecture 8 told us a random projection preserves
the image geometry to within a stated tolerance --- and it did, the pairwise
cosines before and after correlating at 0.86. \emph{It said nothing at all
about aligning that geometry with anything else.}

A similarity is meaningless alone, so compute two numbers: the cosine between
an image and its own caption, and between an image and somebody else's ---
nearly forty thousand of the latter.

\begin{center}
\small
\begin{tabular}{lrr}
\toprule
\textbf{Pairs} & \textbf{Mean cosine} & \textbf{sd} \\
\midrule
Image and its own caption & $+0.0060$ & 0.0499 \\
Image and an unrelated caption & $+0.0034$ & 0.0514 \\
\midrule
Difference & $+0.0026$ & Cohen's $d = 0.052$ \\
\bottomrule
\end{tabular}
\end{center}

\begin{center}
\small
\begin{tabular}{lrrrr}
\toprule
\textbf{System} & \textbf{R@1} & \textbf{R@5} & \textbf{R@10} &
\textbf{Median rank} \\
\midrule
Random ranking (exact) & 0.5\% & 2.5\% & 5.0\% & 100 \\
ViT + MiniLM, JL projection & 1.0\% & 4.0\% & 6.0\% & 104 \\
ViT + MiniLM, truncated & 0.5\% & 2.0\% & 5.5\% & 80 \\
ViT + MiniLM, zero-padded & 0.0\% & 2.5\% & 5.5\% & 80 \\
\bottomrule
\end{tabular}
\end{center}

Two of the 200 queries put the right image first. Shuffling is expected to get
one. Two strong encoders, over a hundred million parameters between them, and
that is the whole difference.

\begin{keybox}[title={Say the result out loud}]
A picture of a bus and the sentence ``a red bus on a city street'' are no
closer than that picture and a sentence about a sandwich.

\emph{Both encoders work perfectly.} The ViT separates images from one another;
MiniLM separates sentences from one another. Neither was ever shown a single
(image, sentence) pair, so neither has any reason to place them anywhere in
particular relative to each other.
\end{keybox}

Why, precisely: an embedding space is defined only up to whatever the training
objective constrained. ImageNet classification constrains images relative to
images; sentence-similarity training constrains sentences relative to
sentences; and \emph{nothing in either loss ever compared a vector from one to
a vector from the other}. Apply any orthogonal map to the image space and every
image--image cosine is unchanged, so the ViT's loss is unchanged --- while
every image--text cosine changes. The loss cannot distinguish those worlds, so
it did not pick one.

\section{Derivation: the contrastive objective and its temperature}

Two encoders trained apart gave 1.0\%. Two trained together give 73.0\%. What
is the loss that did that? Four decisions: where the vectors live, what an
unrelated pair should score, how a cosine becomes a logit, and how many wrong
answers the loss compares against.

A batch of $B$ pairs, each tower ending in a linear map into $\mathbb{R}^d$,
gives a $B\times B$ matrix of cosines $S_{ij}$. $S_{ii}$ is a pair that belongs
together; every off-diagonal entry is a pair that does not. \emph{The loss only
ever sees this matrix.}

\subsection{Why normalise}

The unnormalised inner product factorises as
$\lVert a\rVert\lVert b\rVert\cos\theta$, mixing two quantities: which
direction the encoder chose, and how loudly it said it. Only the direction
carries the semantics --- the length is whatever the last linear layer happened
to scale to, and on our catalogue the longest image embedding is 1.32 times the
shortest.

Normalising costs one degree of freedom per vector and the ability to express
confidence by shouting; it buys a bounded score, so that a single temperature
is meaningful for every pair, and a geometry we can reason about. Confidence
comes back through the temperature instead, as one number shared by the whole
model rather than one per example.

\subsection{What should an unrelated pair score?}

The answer everybody gives first is $-1$. It is wrong, and the reason is
geometry already measured in Lecture 8.

For a given $\mathbf{u}$ there is exactly \emph{one} unit vector at cosine
$-1$, namely $-\mathbf{u}$. A catalogue of 200 entries would need all 200
non-matching captions to sit on that one point --- and then they would all be
identical to each other. Wanting every unrelated pair at $-1$ is asking for a
configuration that does not exist for more than two items.

\begin{derivbox}[title={Concentration on the sphere}]
Draw $\mathbf{u}, \mathbf{v}$ uniformly on $S^{d-1}$. Fix $\mathbf{u} = e_1$ by
rotational symmetry and write $\mathbf{v} = \mathbf{z}/\lVert\mathbf{z}\rVert$
with $\mathbf{z} \sim \mathcal{N}(0, I_d)$. Then
$\mathbf{u}^{\top}\mathbf{v} = z_1/\lVert\mathbf{z}\rVert$ and
\[ \mathbb{E}\!\left[\Bigl(\tfrac{z_1}{\lVert\mathbf{z}\rVert}\Bigr)^{2}\right]
   = \mathbb{E}\!\left[\frac{z_1^2}{\sum_k z_k^2}\right] = \frac{1}{d} \]
by symmetry --- the $d$ coordinates share the total equally. The expectation is
zero by $\mathbf{v}\mapsto-\mathbf{v}$, so the standard deviation is exactly
$1/\sqrt d$. No approximation anywhere.
\end{derivbox}

\emph{In high dimensions, two unrelated things are orthogonal, not opposite.}

\begin{center}
\small
\begin{tabular}{lr}
\toprule
\textbf{Random unit vectors in 512 dimensions, 20{,}000 pairs} &
\textbf{Value} \\
\midrule
Mean cosine & $-0.00072$ \\
Standard deviation, measured & 0.0439 \\
Standard deviation, $1/\sqrt d$ & 0.0442 \\
Most negative cosine observed & $-0.177$ \\
Fraction with $|\cos| > 0.5$ & 0.0\% \\
\bottomrule
\end{tabular}
\end{center}

The most negative cosine in twenty thousand draws was $-0.18$. Nothing came
anywhere near $-1$, and nothing ever will. \emph{That is why the loss targets
zero.}

A related consequence worth naming: images occupy one region of the sphere and
captions another, with the two centroids 0.83 apart. The objective aligns pairs
\emph{relatively} and never asks the two clouds to coincide --- the modality
gap survives training.

\subsection{From a cosine to a logit}

Row $i$ of $S$ is a $B$-class problem whose correct answer is column $i$, so
Lecture 17's machinery applies directly:
\[ \mathcal{L} = -\frac{1}{B}\sum_i
   \log\frac{\exp(S_{ii}/\tau)}{\sum_j \exp(S_{ij}/\tau)}. \]

\begin{failbox}[title={Why $\tau = 1$ cannot work}]
The logits are cosines, so the whole spread of a row is at most 2, and
$\exp$ of a range of 2 is a ratio of at most $e^2 \approx 7.4$ across 200
competitors. The softmax is nearly uniform whatever the model says.

Measured on our 200 pairs at $\tau = 1$: loss 5.151, against 5.298 for a scorer
that knows nothing, with mean probability 0.0058 on the correct entry.
\end{failbox}

Differentiating gives Lecture 17's result scaled:
$\partial\mathcal{L}_i/\partial S_{ij} = \frac{1}{\tau}(p_{ij} - \mathbf{1}[j
= i])$. So $\tau$ does not change which column is largest --- \emph{it cannot
change the top-1 accuracy of a fixed model} --- it changes where the gradient
goes. Small $\tau$ concentrates the weight on the few hardest negatives.

\begin{center}
\small
\begin{tabular}{lrrr}
\toprule
$\tau$ & \textbf{Loss} & \textbf{p(correct)} & \textbf{p(hardest wrong)} \\
\midrule
1 & 5.151 & 0.0058 & 0.0056 \\
0.0100 (learned) & 0.836 & 0.667 & 0.208 \\
0.002 & colder & 0.724 & 0.252 \\
\bottomrule
\end{tabular}
\end{center}

Colder is not simply better: as $\tau \to 0$ the softmax becomes a hard
maximum, the gradient sees one negative per row, and any label noise in that
one pair is amplified rather than averaged away.

And $\tau$ is not a hyperparameter anybody tunes by hand --- the model stores
$\log(1/\tau)$ and learns it, clamped from above. Read off the checkpoint, the
learned scale is 100.00 to five significant figures. \emph{It is not near its
clamp; it is on it}: the implementation caps $\log(1/\tau)$ at $\log 100$, and
training pushed it there and held it.

\subsection{Where the negatives come from}

Nowhere --- that is the elegant part. The loss needs, for each image, a set of
captions it should not match, and nobody labels those: they are simply the
other members of the batch, free and correct with high probability on a large
corpus.

\begin{keybox}[title={The batch is the label set}]
So the batch size is not a memory setting. It is \emph{the number of classes in
the classification problem you are solving.}
\end{keybox}

\begin{center}
\small
\begin{tabular}{rrrr}
\toprule
\textbf{Batch size $B$} & \textbf{In-batch top-1} & \textbf{Chance, $1/B$} &
\textbf{Loss} \\
\midrule
  2 & 99.2\% & 50.0\% & 0.016 \\
 16 & 95.8\% &  6.2\% & 0.116 \\
 64 & 87.4\% &  1.6\% & 0.389 \\
200 & 72.5\% &  0.5\% & 0.836 \\
\bottomrule
\end{tabular}
\end{center}

Same model, same embeddings, same temperature --- only the number of
distractors changed. Accuracy falls with $B$ and chance falls faster, so the
gap widens: that is the whole argument for a large batch.

\begin{failbox}[title={Two consequences worth carrying away}]
A contrastive loss value is \emph{not comparable across papers}, because it is
measured against a batch-dependent ceiling of $\log B$. Our own sweep makes the
point: the same model, embeddings and $\tau$ score 0.058 at $B = 8$ and 0.836
at $B = 200$. Neither number means anything without its batch size.

And batch size stops being an optimisation detail. Doubling it changes the
\emph{task}, not just the gradient noise --- which is why these models are
trained with batches in the tens of thousands.
\end{failbox}

Lecture 8 supplied the concentration result; Lecture 17 supplied the softmax
and its gradient. \emph{The eighteen derivations were meant to be taken in
order, and this is why.}

\section{One space, two modalities}

A jointly trained dual encoder --- a ViT-B/32 image tower, a twelve-layer text
tower, and a shared 512-dimensional space on the unit sphere. The whole system
is fourteen lines, and the only difference from the failed attempt is that the
two sides now have the same shape, so the matrix product exists.

\begin{failbox}[title={One line worth staring at}]
\code{get\_image\_features} and \code{get\_text\_features} return vectors that
are \emph{not} normalised. Skip the normalisation and the code still runs,
still returns a matrix, and still produces a ranking --- a different one.
\end{failbox}

\subsection{The known-answer test first}

The catalogue has no labels, so nothing there can tell us the model is loaded
correctly, pointed at the right tensors, and preprocessing images the way it
was trained to. So: 2{,}000 CIFAR-10 images, ten class names, embed ``a photo
of a \{class\}'' for each, and assign each image to the nearest sentence.
\emph{The classifier is ten sentences, and there is no fitted parameter
anywhere.}

90.5\% over 2{,}000 images against a 10.0\% chance level, from a model that has
never seen CIFAR-10.

\begin{keybox}[title={The sentence is a hyperparameter}]
\begin{center}
\small
\begin{tabular}{lr}
\toprule
\textbf{Prompt} & \textbf{Accuracy over 2{,}000 images} \\
\midrule
``dog'' & 88.2\% \\
``a photo of a dog'' & 90.5\% \\
``a low-resolution photo of a dog'' & 91.0\% \\
\bottomrule
\end{tabular}
\end{center}
A spread of 2.8\% from changing the wording alone. \emph{A ``zero-shot'' number
is a number for one particular sentence, and the sentence is a choice you
made.}
\end{keybox}

\subsection{Back to the catalogue}

Same 200 images, same 200 queries, same metric, same ties rule, same baseline
--- only the encoders changed. In the joint space the matched and unrelated
distributions come apart: the means differ by $+0.149$, which is 4.1 pooled
standard deviations against the $+0.0026$ and $d = 0.052$ of the separate
encoders.

\begin{center}
\small
\begin{tabular}{lrrrrr}
\toprule
\textbf{System} & \textbf{R@1} & \textbf{R@5} & \textbf{R@10} &
\textbf{Median rank} & \textbf{MRR} \\
\midrule
Random ranking & 0.5\% & 2.5\% & 5.0\% & 100 & 0.029 \\
ViT + MiniLM & 1.0\% & 4.0\% & 6.0\% & 104 & 0.037 \\
Joint dual encoder & 73.0\% & 96.0\% & 99.5\% & 1 & 0.827 \\
\bottomrule
\end{tabular}
\end{center}

Text to image, 200 candidates, ties against the model, one caption per entry as
the query. \emph{Two of those three rows are the same row.}

\subsection{Report the failure cases}

73.0\% at rank 1 means 27.0\% of queries did not get the right entry first, and
they fall into groups: queries describing a scene rather than an object, where
several entries fit equally well and the ``wrong'' one is often a reasonable
answer; fine distinctions the caption makes and the picture barely shows; and
counts and spatial relations --- ``three'', ``behind'', ``to the left of''.

\begin{keybox}[title={One relevant entry is a fiction}]
Our metric assumes exactly one right answer per query because that is what the
data gives us. For ``a man riding a horse'' on a real catalogue it is simply
false, so the measured recall is a \emph{lower bound}.

And customers do not type captions. They type ``something to sit on'', ``gear
for bad weather''. The right output there is not one entry but a shortlist with
reasons --- and we have no way to produce, or to score, a reason.
\end{keybox}

\subsection{Six ways to leak in a retrieval evaluation}

The query text is also the indexed description, so you are measuring string
equality. The query's own entry is scored against a shortlist built using the
query. The prompt template was chosen by looking at the number it produced.
Recall is reported without the candidate-set size. Ties are broken in the
model's favour. Or normalisation is applied to one side only, so ``cosine'' is
not a cosine.

\begin{failbox}[title={The tie bug, which reports 100\%}]
Ranking with a strict inequality counts only candidates that \emph{beat} the
right answer, so anything tying with it is placed behind it for free.

Feed that a similarity matrix of all zeros --- a model that has learned
nothing, or a bug that zeroed the embeddings --- and every rank is 1. It
reports R@1 = 100\%.
\end{failbox}

And the third leak above is Lecture 2's, with a sentence in place of a
regularisation strength: three templates tried, the best on the test set
reported. The spread across the three was 2.8\%, so the optimism is of that
order, and the repair is the one you already know.

\section{The notebook}

\begin{trybox}[title={What to change}]
Skip the normalisation before the inner product. Retrieval degrades, and the
reason is in the derivation: without unit norm the score confounds direction
with magnitude, and magnitude carries no meaning here.
\end{trybox}

\section{Exercises}

\begin{enumerate}[leftmargin=1.6em,itemsep=4pt]
  \item Two encoders are trained separately, one on images and one on text.
        Explain why their cosine similarity carries no information about
        matching, using a rotation argument. \hfill \scope{[5]}
  \item In $d$ dimensions two independent random unit vectors have a cosine
        with standard deviation $1/\sqrt d$. State what an unrelated pair looks
        like at $d = 512$, and why $-1$ is the wrong intuition.
        \hfill \scope{[5]}
  \item A contrastive loss is computed at $\tau = 1$ on cosine scores. State
        why it barely trains, and what $\tau$ changes and does not change.
        \hfill \scope{[5]}
  \item Recall@10 is reported over 200 candidates. State the exact random
        baseline, and why it is computed rather than simulated.
        \hfill \scope{[3]}
  \item An entry with no written description scores zero on a text retrieval
        route. State whether this is a ranking failure or a structural one, and
        what follows. \hfill \scope{[4]}
\end{enumerate}

\section*{Summary}
\addcontentsline{toc}{section}{Summary}

An image becomes a sequence of patches, which gives up exactly what Lecture 12
derived convolution assumes --- locality and translation equivariance --- and
buys the ability to learn whichever relations actually hold, at a corpus size
most people do not have.

Two separately trained encoders produce spaces that are \emph{unrelated}: the
matched and unmatched cosine distributions differ by 0.0026 at a Cohen's $d$ of
0.052, and the retrieval is indistinguishable from shuffling. Not because
either encoder is bad, but because no loss ever compared a vector from one with
a vector from the other, and an orthogonal map on the image space leaves that
loss unchanged while changing every image--text cosine.

The contrastive objective fixes four things. Normalise, because only direction
is semantic. Target zero rather than $-1$ for unrelated pairs, because on the unit
sphere in 512 dimensions random directions have cosine spread $1/\sqrt d$ --- predicted
0.0442, measured 0.0439, with the most negative of 20{,}000 draws at $-0.18$.
Divide by a learned $\tau$, because cosines in $[-1,1]$ make a softmax nearly
uniform; it cannot change a fixed model's ranking, only which negatives the
gradient hears. And take the negatives from the batch, so $B$ sets both the
difficulty and the chance level --- which makes a contrastive loss value
meaningless without its batch size.

Jointly trained, the same measurement gives R@1 of 73.0\% against 1.0\%, and
zero-shot classification reaches 90.5\% on a corpus the model has never seen
--- for one particular prompt, out of three that span 2.8\%.

\end{document}
